\chapter{来源与史料说明}
\label{app:sources}

本附录是本卷材料的阅读索引，不是一张把所有文献排成高低次序的清单。周秦的材料保存很不均匀：早周可直接依凭的金文和考古遗存很珍贵，却往往只留下一个受命、一次赏赐或一段宗族记忆；春秋、战国和秦则有更多编年与后出的叙事，仍不能因此把它们当作逐日的档案。阅读每一事件时，先问它是哪一类证据、最接近什么问题、又不能回答什么问题。

\section{按时段进入材料}

\begin{tabularx}{\linewidth}{@{}p{0.17\linewidth}p{0.34\linewidth}X@{}}
\toprule
\textbf{时段} & \textbf{优先材料} & \textbf{可用范围与边界} \\
\midrule
西周 &
\sourcebook{尚书}若干篇、\sourcebook{诗经}、\sourcebook{史记·周本纪}、\sourcebook{竹书纪年}，以及利簋、何尊、大盂鼎等金文和城址、墓葬材料 &
传世文本保存王朝记忆、政治语言和后出的编年整理；金文可旁证称号、王命、赏赐、地名和精英网络。两者都不能自动复原某次战役的部署、封国的精确边界或普通人的完整生活账目。 \\
\addlinespace
春秋 &
\sourcebook{春秋}、\sourcebook{左传}、\sourcebook{国语}、\sourcebook{公羊传}、\sourcebook{谷梁传}、\sourcebook{史记}相关世家，并参照盟书、青铜器、城址和兵器 &
\sourcebook{春秋}提供鲁国年次和行动的最小骨架；传世叙事使外交、战争和礼仪关系可被讨论，却不是会盟记录或战场速记。区域考古能检验空间与物质条件，不能替代失传的谈判和军政档案。 \\
\addlinespace
战国与秦 &
\sourcebook{史记}本纪、世家与列传，\sourcebook{资治通鉴}，\sourcebook{战国策}，\sourcebook{商君书}、\sourcebook{韩非子}等思想材料，以及睡虎地、里耶秦简、秦刻石、封泥、权量器和遗址 &
后出的史书与编年著作能保持国家线和年代脊柱；说辞和思想文本更适合观察后人如何组织政治语言，而非逐字还原游说或法令。简牍、刻石和器物可见局部行政实践，不能由一地、一时的文书推出全国同步运行。 \\
\bottomrule
\end{tabularx}

\section{按问题交叉核对}

\begin{description}
  \item[年代与政治行动] 先以最接近该事件的编年线索确定最低限度的顺序，再以国别叙事和后出的综合史书比较异文。牧野、共和、春秋战国分期与秦亡的边界尤其不可只凭一个整齐的年代表判断；可从\eventref{WZ-1046-MUYE}{西周·牧野}、\eventref{WZ-0841-GONGHE}{西周·共和}、\eventref{SA-0481-END}{春秋·记事终点}和\eventref{QIN-0209-CHENSHENG}{秦·大泽起事}返回正文与争议说明。
  \item[战争、外交与地理] 编年和国别叙事说明哪些政治体参与、冲突如何被记忆；河流、道路、城址、兵器和区域调查可检验行动的空间条件。它们可以帮助读者重读\eventref{SA-0632-CHENGPU}{春秋·城濮之战}、\eventref{SA-0506-BOJU}{春秋·柏举之战}与\eventref{WQ-0260-CHANGPING}{战国·长平之战}，却不能单独裁定谋略、兵数或人物动机。
  \item[制度、征发与地方社会] 诰命、金文、法令性文本、简牍和后出的制度叙述各自只照亮一部分问题。阅读\eventref{WZ-EAST-ENFEOFFMENT}{西周·重组东方封国网络}、\eventref{WQ-0356-SHANGYANG}{战国·商鞅变法}与\eventref{QIN-0213-EMPIRE}{秦·高强度治理}时，应分开制度理想、个案文书、后世概括和实际执行，尤其不要把秦汉以后的成熟制度倒投回更早的时代。
  \item[文本、思想与接受] \sourcebook{论语}、\sourcebook{墨子}、\sourcebook{孟子}、\sourcebook{庄子}、\sourcebook{荀子}、\sourcebook{吕氏春秋}及法家文本，是讨论概念、论辩与后来接受史的重要入口；其成书层次、传本和人物归属本身也是问题。它们不能替代政治事件的独立证据。
\end{description}

\section{本卷的来源路径}

\begin{description}
  \item[先秦传世叙事] \sourcebook{尚书}、\sourcebook{诗经}、\sourcebook{春秋}与三传、\sourcebook{左传}、\sourcebook{国语}、\sourcebook{战国策}和\sourcebook{史记}构成本卷最常用的文字材料群。它们的文体不同：诰命、诗篇、经年、叙事、说辞和汉代史书不能互相替代。
  \item[编年参照] \sourcebook{竹书纪年}提供异年表线索；\sourcebook{资治通鉴}将战国至秦末重新编年，便于追踪前后关系。两者都须与更接近具体事件的材料及现代考订互校。
  \item[出土与物质材料] 金文、盟书、简牍、封泥、刻石、权量器、城址、墓葬与战场遗存，优先用于称谓、文书、尺度、空间、技术和精英网络。出土背景、断代和释读不同，材料能够承担的结论也不同；无可靠背景的器物或未核实释文不作为本卷的核心依据。
  \item[现代研究] 年代学、古文字与文本校勘、考古报告、历史地理、战争与交通研究、制度史和社会经济史，分别用来检验材料的年代、传本、空间条件与解释尺度。吕思勉、钱穆等总体史著作可帮助提出连续性的提问，不替代逐年事实的核对。
\end{description}

\section{基础书目与读法索引}
\label{app:bibliography}

下列条目只收录本卷来源路径已经点明的材料，以及书系规范要求读者辨明用途的解释性读物。现有来源矩阵没有保存出版社、版次、页码或逐条引文的版本记录，故本附录不虚构一个“本卷所用的唯一版本”：每一条均明确标为\textbf{本卷未锁定的版本状态}。在编辑补入可复核的整理者、出版社与版次前，这里是读者的检索与读法入口，不是可替代正文证据链的书目引文。

\subsection{本卷材料的检索入口}

\begin{description}
  \phantomsection\label{src:western-texts}
  \item[\sourcebook{尚书}、\sourcebook{诗经}] 作者/编者与成书层次：传世篇章与诗篇的汇集，具体篇章的年代、传本与作者归属并不整齐；版本：本卷未锁定整理本。适用的是早周政治语言、受命记忆与礼制表达；本卷借此说明传统怎样塑造周初叙事，不能把篇章直接还原为某日诰命、战报或制度清单。

  \phantomsection\label{src:chunqiu}
  \item[\sourcebook{春秋}] 作者/编者与成书层次：传世鲁国编年传统；版本：本卷未锁定整理本。它为前722至前481年的鲁国年次、会盟、伐战与继承提供最小骨架；它沉默之处不能倒推出列国没有行动，也不能单靠极简经文说明动机和战术。

  \phantomsection\label{src:zuozhuan-guoyu}
  \item[\sourcebook{左传}、\sourcebook{国语}、\sourcebook{公羊传}、\sourcebook{谷梁传}] 作者/编者与成书层次：传统归属与实际成书、编纂层次均有争议；版本：本卷未锁定整理本。前二者用于比较国别政治记忆、外交与礼仪叙事，后二者用于理解经学解释如何赋义；它们都不能把说辞、预言或褒贬义例变成会盟实录或独立的事件证据。

  \phantomsection\label{src:shiji}
  \item[\sourcebook{史记}] 作者：司马迁；版本：本卷未锁定点校本。它是本卷连接周、列国与秦亡的核心后出叙事，用于国别线索、异文与汉代如何组织早期历史；不能替代同代金文、简牍或失传的军政档案，也不能将人物传记的戏剧性细节当作现场记录。

  \phantomsection\label{src:bamboo-annals-zizhi}
  \item[\sourcebook{竹书纪年}、\sourcebook{资治通鉴}] 编者：前者的传本、辑佚与重建关系复杂，后者由司马光主持编纂；版本：本卷均未锁定整理本。前者提供异年表线索，后者帮助追踪战国至秦末的前后关系；二者都须与更接近具体事件的材料互校，不能单独裁定西周年代或替代周秦的国别材料。

  \phantomsection\label{src:warring-texts}
  \item[\sourcebook{战国策}、\sourcebook{商君书}、\sourcebook{韩非子}、\sourcebook{管子}、\sourcebook{吕氏春秋}] 编者/作者：\sourcebook{战国策}为刘向整理的国策说辞传统；其余著作虽有传统作者或相关政治人物，篇章成书与传本层次并不一致；版本：本卷未锁定整理本。它们用于考察游说、制度语言和后人怎样理解改革；不能逐字还原游说、诏令或某位人物在某年的个人主张。

  \phantomsection\label{src:thought-texts}
  \item[\sourcebook{论语}、\sourcebook{墨子}、\sourcebook{孟子}、\sourcebook{庄子}、\sourcebook{荀子}、\sourcebook{孙子兵法}] 作者/编者与成书层次：均有传本、篇章或归属问题；版本：本卷未锁定整理本。它们是理解先秦思想、论辩与后世接受的阅读入口；并非政治事件的独立编年证据，不能以一段名言替代年代、政策或人物行动的核对。

  \phantomsection\label{src:excavated}
  \item[金文、盟书、简牍与考古报告] 编者/版本：利簋、何尊、大盂鼎、侯马盟书、睡虎地与里耶秦简及相关城址、墓葬、兵器、封泥、刻石、权量器，是按器物、发掘报告和释读分别检索的材料群；本卷来源矩阵未记录可逐项复核的报告版本。它们用于称谓、王命、文书、地方行政、空间与物质条件的旁证；一地一时或一件器物不得外推为全国同步制度，更不能补写传世叙事未保存的谈判和战场细节。
\end{description}

\subsection{解释性读物与古文阅读}

\begin{description}
  \phantomsection\label{src:interpretive}
  \item[钱穆《国史大纲》《历代政治得失》；吕思勉通史与断代著作；张荫麟《中国史纲》] 作者如题；版本：本卷未锁定版本，亦未将它们作为逐条事件的独立引文。它们可帮助读者提出时代大势、制度得失、问题意识以及社会经济长线的问题；不能单独证明逐年事实、精确数字或人物动机。若其概括与原始材料或考古研究不合，应回到具体材料，而非以通史概括裁决。

  \phantomsection\label{src:guwen-guanzhi}
  \item[\sourcebook{古文观止}] 编者：吴楚材、吴调侯；版本：本卷未锁定版本，也未把它作为事件证据。它可以是阅读古文的选本入口；需要使用其中篇章时，仍应回到原始文本、可靠底本与成文语境，不能把选本篇名或注释当作原始史料。
\end{description}

\subsection{书系正史骨架}

这一表是书系的检索骨架，不声称这些史书都已被第一卷逐页使用；其中只有\sourcebook{史记}是本卷实际的核心叙事来源之一。各卷应在自己的来源附录锁定所用整理本和事件路径，不能把下表的覆盖提示误当作本卷的版本引用。

\begin{description}
  \item[第一卷·周秦] \sourcebook{史记}：司马迁撰；用于周秦的后出叙事骨架，限度见\hyperref[src:shiji]{本条目}。
  \item[第二卷·汉帝国] \sourcebook{汉书}、\sourcebook{后汉书}：两汉断代史的主体；仍须与编年、制度志、文书和出土材料互校。
  \item[第三卷·三国] \sourcebook{三国志}：区域政权与人物线索的主体；裴松之注及后出叙事应与本文区分。
  \item[第四卷·两晋南北朝] \sourcebook{晋书}、\sourcebook{宋书}、\sourcebook{南齐书}、\sourcebook{梁书}、\sourcebook{陈书}、\sourcebook{魏书}、\sourcebook{北齐书}、\sourcebook{周书}、\sourcebook{南史}、\sourcebook{北史}：须同时面对修史年代、政权立场和南北材料不均。
  \item[第五卷·隋唐五代] \sourcebook{隋书}、\sourcebook{旧唐书}、\sourcebook{新唐书}、\sourcebook{旧五代史}、\sourcebook{新五代史}：须与诏令、墓志、地方与专门编年材料并读。
  \item[第六卷·宋辽夏金元] \sourcebook{宋史}、\sourcebook{辽史}、\sourcebook{金史}、\sourcebook{元史}：须与\sourcebook{续资治通鉴长编}、\sourcebook{建炎以来系年要录}、\sourcebook{宋会要辑稿}及多政权文书、碑刻和考古互校。
  \item[第七卷·明] \sourcebook{明史}：须与\sourcebook{明实录}、档案、地方文献和物质材料交叉，不能把后出的总纂叙事当作唯一记录。
  \item[第八卷·清至1912] 清代不在二十四史内；以\sourcebook{清实录}、\sourcebook{清史稿}、档案、地方文献与近代报刊等接续。它们各有官修、保存与立场边界，不能互相替代。
\end{description}

\section{怎样沿索引复核}

读者可先从\hyperref[app:index]{“周秦卷索引”}按事件、国家、主题、来源或图件进入正文，再回到本附录判断所用材料的类型；需要了解哪些不确定性会改变叙事结论时，参看\hyperref[app:controversies]{“争议怎样改变周秦史的读法”}；需要由材料名称进入读法与版本状态时，参看\hyperref[app:bibliography]{“基础书目与读法索引”}。本卷不以引文数量制造确定性：一条来源链的价值，在于它让读者看见何处是可确认的行动，何处是传世叙事的组织，何处仍须保留空白。
